\section{Desglose de un caso: cómo diseñar un experimento de mejora de reasoning}

\subsection{Objetivos y restricciones del experimento}
Supongamos que queremos mejorar un modelo de la escala 7B en 5 a 10 puntos en tareas de razonamiento matemático y de código, sin dañar de manera notable la capacidad de conversación general. La restricción de presupuesto es ``entrenamiento controlable, inferencia lista para producción''. Este tipo de problema corresponde exactamente a la ruta combinada que propone esta clase.

\begin{importantbox}{Cómo definir los objetivos}
El objetivo debe incluir al mismo tiempo:
\begin{itemize}
    \item \textbf{Objetivo de capacidad}: mejora de la exactitud en tareas verificables;
    \item \textbf{Límite de efectos secundarios}: la caída en tareas generales no debe superar un umbral;
    \item \textbf{Objetivo de despliegue}: costo y latencia dentro de un rango aceptable.
\end{itemize}
Si falta cualquiera de estos elementos, la dirección de la optimización se desvía de la necesidad real.
\end{importantbox}

\subsection{Matriz de experimentos por etapas}
\begin{table}[H]
\centering
\renewcommand{\arraystretch}{1.2}
\begin{tabular}{p{2.8cm}p{3.4cm}p{3.8cm}p{3.0cm}}
\toprule
\textbf{Etapa} & \textbf{Variable} & \textbf{Métrica observada} & \textbf{Condición de detención} \\
\midrule
SFT & Proporción de datos humanos/sintéticos, distribución de longitud & Seguimiento de instrucciones, estabilidad de formato, exactitud básica & Revertir en cuanto el efecto secundario supere el umbral \\
DPO/PPO & Peso de la pérdida de preferencia, temperatura de muestreo & Utilidad, verbosidad, tasa de rechazo de respuesta & El beneficio marginal de la preferencia disminuye \\
RLVR & Reglas del verificador, número de rondas de muestreo, número de pasos de actualización & Tasa de aprobación verificable en matemáticas/código & Aparece una señal de atajo de recompensa \\
Inference-time & Valor de N, estrategia de reordenamiento, número de rondas de recuperación & Tasa de aprobación, latencia P95, costo por solicitud & El costo supera el presupuesto \\
\bottomrule
\end{tabular}
\caption{Matriz por etapas de un experimento de mejora de reasoning}
\end{table}

\subsection{Prioridad del análisis de errores}
\begin{knowledgebox}{Los tres tipos de muestras de error más valiosos}
\begin{enumerate}
    \item \textbf{Error de alta confianza}: el modelo produce un razonamiento largo y encadenado, expresado con confianza pero incorrecto;
    \item \textbf{Juicio erróneo del verificador}: la respuesta es correcta pero la regla la descarta por error, o al revés;
    \item \textbf{Muestras de presupuesto ineficiente}: requieren un muestreo masivo para lograr una mejora mínima en la exactitud.
\end{enumerate}
Corregir primero estas muestras suele mejorar al mismo tiempo la calidad y la eficiencia de costos.
\end{knowledgebox}

\begin{warningbox}{No ocultar fallas estructurales con el puntaje promedio}
Un aumento en el puntaje promedio puede ocultar el retroceso en tareas clave. Un sistema real necesita examinar más bien ``el desempeño por segmento de tareas'': por ejemplo, cómo cambian por separado las matemáticas de varios pasos, la reparación de código largo y las preguntas y respuestas que dependen de recuperación.
\end{warningbox}

\subsection{Cómo relacionar este caso con el hilo principal de la clase}
Este caso muestra que el marco de tres etapas y la idea de receta abierta de esta clase pueden traducirse directamente en decisiones de ingeniería:
\begin{itemize}
    \item El preentrenamiento determina el límite superior del modelo base;
    \item El posentrenamiento determina la dirección del comportamiento;
    \item El escalamiento en tiempo de inferencia determina la curva de beneficio en el lado del despliegue;
    \item La apertura y la reproducibilidad determinan si el equipo puede iterar de manera sostenida.
\end{itemize}

\subsection{Resumen de la sección}
Un experimento exitoso de mejora de reasoning no es la victoria de un solo parámetro puntual, sino el resultado de que las restricciones del objetivo, la matriz de etapas, el análisis de errores y la gestión del presupuesto trabajen de manera coordinada.

